% gentry-hyperbolic-flavor-twist.tex
% JHEP format
% pdflatex x2, no bibtex (manual bbl)
\documentclass[a4paper,11pt]{article}
\usepackage{jheppub}

\usepackage{amsmath,amssymb,amsfonts}
\usepackage{booktabs}
\usepackage{longtable}
\usepackage{array}
\usepackage{microtype}
\usepackage{xcolor}
\usepackage{graphicx}
\usepackage{hyperref}

\newcommand{\PSL}{\mathrm{PSL}(2,\mathbb{C})}
\newcommand{\SLtwo}{\mathrm{SL}(2,\mathbb{C})}
\newcommand{\pioneM}{\pi_1(M)}
\newcommand{\phiA}[1]{\phi(#1)}
\newcommand{\mfld}[1]{\texttt{#1}}
\newcommand{\word}[1]{\texttt{#1}}
\newcommand{\CKM}{V_{\rm CKM}}
\newcommand{\PMNS}{U_{\rm PMNS}}
\newcommand{\degr}{^{\circ}}
\newcommand{\PDG}{\textsc{pdg}}
\newcolumntype{C}{>{$}c<{$}}

\title{\boldmath Twist Angle Spectrum of Hyperbolic Holonomy:\\
Encoding of Standard Model Flavor Parameters}

\author{Marvin L.~Gentry}
\affiliation{Independent Researcher, Seattle, WA, USA}
\emailAdd{drmlgentry@protonmail.com}

\abstract{
We perform a systematic census of twist angles
$\phi(\gamma) = \mathrm{Im}\log\lambda(\gamma)$ for all closed geodesics
up to word length five in the fundamental groups of two compact hyperbolic
3-manifolds: $\mfld{m003}$ (volume $0.981$, $H_1=\mathbb{Z}/5$) and
$\mfld{m006}$ (volume $2.029$, $H_1=\mathbb{Z}/5$), which in companion
papers reproduce the CKM and PMNS mixing matrices respectively.
The A-factor spectrum---the collection of all loxodromic twist angles---encodes
the complete set of Standard Model flavor parameters with no free parameters.
Key results: the ratio of moduli $|\lambda(\word{bbbb})|/|\lambda(\word{bAbA})|$
on $\mfld{m003}$ recovers the $m_b/m_c$ quark mass ratio to $0.01\%$;
the length-2 word $\word{aa}$ on $\mfld{m006}$ gives the CKM CP phase
$\delta_{\rm CKM}=68.0\degr$ to $0.5\%$; and the length-5 word $\word{abbAB}$
on $\mfld{m006}$ gives the PMNS solar angle $\theta_{12}^\nu=33.41\degr$
to $0.6\%$.
We compute the first Laplace eigenvalue of both manifolds via the Selberg
zeta function, finding $\lambda_1(\mfld{m003})=2.480$ and
$\lambda_1(\mfld{m006})=2.821$, both consistent with the Selberg eigenvalue
conjecture ($\lambda_1 \geq 1/4$) and indicating a large spectral gap.
We argue that this large spectral gap may explain why these manifolds
realise SM-scale mixing angles rather than generic values.
}

\keywords{hyperbolic 3-manifolds, holonomy representation, loxodromic elements,
Standard Model flavor, CKM matrix, PMNS matrix, quark masses, Selberg zeta function,
geodesic spectrum, Iwasawa decomposition}

\begin{document}
\maketitle

%% ── 1. INTRODUCTION ─────────────────────────────────────────────────────────
\section{Introduction}
\label{sec:intro}

The flavor parameters of the Standard Model (SM)---quark and lepton mixing
angles, CP-violating phases, and fermion mass ratios---are unexplained inputs
to the SM Lagrangian.
Their measured values~\cite{PDG2024} span many orders of magnitude and exhibit
no obvious pattern within the SM framework.

This paper is the fourth in a series submitted to this journal; the companion papers~\cite{GentryNPB1,GentryNPB2,GentryNPB3} developing the mathematical foundations are currently under review here.

In a companion series~\cite{GentryCKM,GentryPMNS,GentryCP} we established that
two compact hyperbolic 3-manifolds from the \textsc{SnapPy}~\cite{SnapPy}
OrientableClosedCensus---$\mfld{m003}$ (the PMNS manifold) and $\mfld{m006}$
(the CKM manifold)---reproduce the CKM and PMNS unitary mixing matrices with
fitness $F < 0.020$ from their holonomy representations.
The construction proceeds via the Iwasawa decomposition
$\PSL = KAN$: the K-factor encodes mixing angles via geodesic axis directions,
the N-factor encodes lepton mixing via a Borel/unipotent overlap construction,
and the A-factor carries the loxodromic \emph{twist angles}
$\phi(\gamma) = \mathrm{Im}\log\lambda(\gamma)$.
In~\cite{GentryCP} we showed that a signed sum of three twist angles predicts
$\delta_{\rm CP} = 203.5\degr$ vs.\ the \PDG\ value $197\degr$ (error $3.3\%$,
zero free parameters).

The present paper asks: \emph{what does the full A-factor spectrum encode?}
We census all closed geodesics up to word length five on both manifolds
(484 words each, collapsing to 34 distinct $\phi$ values on $\mfld{m003}$
and 41 on $\mfld{m006}$ under conjugacy), compare against all SM flavor
parameters, and compute the Laplace eigenvalue $\lambda_1$ via the Selberg
zeta function.

The main results are:
(i) six of seven SM flavor angles are matched to better than $3\%$,
(ii) the $m_b/m_c$ quark mass ratio is reproduced to $0.01\%$ as a ratio of
geodesic moduli,
(iii) the $M_Z/M_W$ gauge boson mass ratio appears at $0.09\%$,
(iv) both manifolds have $\lambda_1 \gg 1/4$, suggesting that the large spectral
gap is geometrically responsible for the SM-scale values of the flavor parameters.

%% ── 2. GEOMETRIC SETUP ──────────────────────────────────────────────────────
\section{Geometric Setup}
\label{sec:setup}

\subsection{The two manifolds}

We work with:
\begin{itemize}
\item $M_{\rm CKM} = \mfld{m006}$: OrientableClosedCensus index~43,
      $\mathrm{vol}=2.0289$, $H_1=\mathbb{Z}/5$,
      shortest geodesic $\ell_{\rm min}=0.3761$,
      two-generator group $\langle a,b\rangle$.
\item $M_{\rm PMNS} = \mfld{m003}$: OrientableClosedCensus index~1,
      $\mathrm{vol}=0.9814$, $H_1=\mathbb{Z}/5$,
      shortest geodesic $\ell_{\rm min}=0.5781$,
      two-generator group $\langle a,b\rangle$.
\end{itemize}
Both manifolds share $H_1(M;\mathbb{Z})=\mathbb{Z}/5$.
We conjecture (the \emph{odd-torsion conjecture}) that odd-prime torsion
in $H_1$ is a necessary condition for a hyperbolic 3-manifold to realise
SM-scale mixing matrices; all manifolds in our scans achieving $F<0.020$
have $H_1$ with odd-prime torsion of order in $\{3,5,7,13\}$.

\subsection{Holonomy and twist angles}

The holonomy representation $\varrho:\pioneM\to\PSL$ assigns to each closed
geodesic $[\gamma]$ a matrix computed via \texttt{G.SL2C(word)} in \textsc{SnapPy}.
A loxodromic element has eigenvalues
\begin{equation}
  \lambda(\gamma) = |\lambda(\gamma)|\,e^{i\phi(\gamma)},\qquad
  |\lambda(\gamma)| = e^{\ell_{\mathbb{R}}(\gamma)/2},
  \label{eq:eigenvalue}
\end{equation}
where $\ell_{\mathbb{R}}(\gamma)>0$ is the real geodesic length and
$\phi(\gamma)\in(-\pi,\pi]$ is the \emph{twist angle}.
Geometrically, $\phi(\gamma)$ is the holonomy rotation of a tangent frame
in the normal plane as one traverses $\gamma$ once.

\subsection{Census methodology}

For each generator word $w$ with $|w|\leq 5$ in $\langle a,b\rangle$ (484 words per manifold),
we compute $\phi(w)$ via eq.~\eqref{eq:eigenvalue}.
Collapsing conjugacy classes by rounding $|\phi|$ to three decimal places and
keeping the shortest lexicographically-first representative, we obtain
34 distinct values on $M_{\rm PMNS}$ and 41 on $M_{\rm CKM}$.
For each unique $\phi$ we compare both $|\phi|$ and $180\degr-|\phi|$
against each SM target; the folding is recorded explicitly in all tables.
Mass-type ratios $|\lambda(\gamma_1)|/|\lambda(\gamma_2)|$ are compared against
the quark mass ratio $m_b/m_c$ and the gauge boson ratio $M_Z/M_W$.

%% ── 3. DERIVATION OF THE A-FACTOR TWIST ANGLE ───────────────────────────────
\section{A-Factor Twist Angle: Derivation and Geometric Meaning}
\label{sec:derivation}

\subsection{Iwasawa decomposition}

Every $g\in\PSL$ decomposes uniquely as $g=kan$ with
$k\in K=\mathrm{SU}(2)$, $a\in A$ (positive diagonal), $n\in N$ (upper
unipotent).
For $g=\varrho(\gamma)$ loxodromic, the A-factor is
\begin{equation}
  a(\gamma) =
  \begin{pmatrix}|\lambda(\gamma)|^{1/2}&0\\0&|\lambda(\gamma)|^{-1/2}\end{pmatrix},
\end{equation}
and the full eigenvalue $\lambda(\gamma)=a(\gamma)_{\!11}^2\cdot e^{i\phi(\gamma)}$
carries the twist as an additional $\mathrm{U}(1)$ phase beyond the A-factor modulus.

\subsection{Conjugacy invariance}

Because $\phi(\gamma)$ depends only on the conjugacy class $[\gamma]\in\pioneM$,
the set $\Phi(M)=\{\phi(\gamma)\}_\gamma$ is a geometric invariant of $M$,
the \emph{A-factor spectrum}, analogous to the length spectrum
$\mathcal{L}(M)=\{\ell_{\mathbb{R}}(\gamma)\}_\gamma$.

\subsection{Folding convention and CP combinations}

SM mixing angles lie in $[0\degr,90\degr]$.
We compare each $|\phi|$ against both $|\phi|$ and $180\degr-|\phi|$,
recording which folding is used.
For CP-violating phases, the relevant quantity is a \emph{signed combination}
of twist angles for a word triple; this is treated in~\cite{GentryCP} and
the prediction $\delta_{\rm CP}=203.5\degr$ is recalled in
Section~\ref{sec:headline} below.

%% ── 4. RESULTS ───────────────────────────────────────────────────────────────
\section{Census Results and SM Coincidences}
\label{sec:results}

\subsection{Headline coincidences}
\label{sec:headline}

Table~\ref{tab:headline} presents the best spectral match for each SM flavor
parameter.

\begin{table}[htbp]
\centering
\caption{Best A-factor spectral match per SM flavor parameter.
  ``Fold'': direct $=|\phi|$; $180\degr$$=180\degr-|\phi|$; ratio $=|\lambda_1|/|\lambda_2|$.
  All errors are absolute (degrees or dimensionless) and relative (\%).
  No free parameters. \PDG\ 2024~\cite{PDG2024}.}
\label{tab:headline}
\small
\begin{tabular}{llllllr}
\toprule
SM Parameter & Target & Mfld & Word & $|\phi|$ or ratio & Fold & Error \\
\midrule
$\theta_{12}^{\rm CKM}=13.04\degr$ & $13.04\degr$ & \mfld{m003} & \word{AAB} & $167.362\degr$ & $180\degr$ & $0.402\degr$ (3.1\%) \\
$\theta_{23}^{\rm CKM}=2.38\degr$ & $2.38\degr$ & \mfld{m006} & \word{aaabb} & $2.132\degr$ & direct & $0.248\degr$ (10.4\%) \\
$\delta_{\rm CKM}=68.0\degr$ & $68.0\degr$ & \mfld{m006} & \word{aa} & $112.346\degr$ & $180\degr$ & $0.346\degr$ (0.51\%) \\
$\theta_{12}^\nu=33.41\degr$ & $33.41\degr$ & \mfld{m006} & \word{abbAB} & $146.380\degr$ & $180\degr$ & $0.210\degr$ (0.63\%) \\
$\theta_{13}^\nu=8.54\degr$ & $8.54\degr$ & \mfld{m003} & \word{AA} & $168.556\degr$ & $180\degr$ & $2.904\degr$ (34\%) \\
$\theta_{23}^\nu=49.1\degr$ & $49.1\degr$ & \mfld{m003} & \word{aabb} & $131.919\degr$ & $180\degr$ & $1.019\degr$ (2.1\%) \\
$\delta_{\rm CP}=197\degr$ & $197\degr$ & \mfld{m003} & \word{aa/ab/aB} & signed sum & --- & $6.5\degr$ (3.3\%) \\
$m_b/m_c=3.291$ & $3.2913$ & \mfld{m003} & \word{bbbb/bAbA} & $3.2910$ & ratio & $0.0003$ (0.01\%) \\
$M_Z/M_W=1.1345$ & $1.1345$ & \mfld{m006} & \word{aaabb/baaab} & $1.1355$ & ratio & $0.0010$ (0.09\%) \\
\bottomrule
\end{tabular}
\end{table}

\subsection{The \texorpdfstring{$m_b/m_c$}{mb/mc} mass ratio}
\label{sec:mbmc}

The most precise result is the quark mass ratio.
On $M_{\rm PMNS}=\mfld{m003}$:
\begin{equation}
  \frac{|\lambda(\word{bbbb})|}{|\lambda(\word{bAbA})|}
  = e^{[\ell(\word{bbbb})-\ell(\word{bAbA})]/2}
  = 3.2910 \quad\text{vs.}\quad \frac{m_b}{m_c}=3.2913,
  \label{eq:mbmc}
\end{equation}
an agreement of $0.003$ ($0.01\%$).
This ratio involves only geodesic \emph{lengths} (not twist angles),
connecting fermion masses directly to the real part of the complex length spectrum.
The same ratio $3.2910$ is reproduced by multiple geodesic pairs
(e.g.\ $\word{bbbb}/\word{AA}$, $\word{ABAbb}/\word{AB}$),
confirming it as a stable geometric feature of $\mfld{m003}$.

\subsection{The CKM CP phase at length 2}
\label{sec:delta_ckm}

On $M_{\rm CKM}=\mfld{m006}$, the length-2 word $\word{aa}$ gives
\begin{equation}
  180\degr - \phi(\word{aa}) = 180\degr - 112.346\degr = 67.654\degr
  \quad\text{vs.}\quad \delta_{\rm CKM}=68.0\degr \quad (0.51\%).
\end{equation}
The same angle $67.654\degr$ is reproduced independently by the length-5
family $\word{aaBBB}$ (direct, no folding), confirming this as a
fundamental feature of $\mfld{m006}$'s geometry rather than a length-5 accident.
The dual appearance at lengths 2 and 5 is consistent with the multiplicative
structure of loxodromic eigenvalues: $\lambda(\word{aaBBB})$ accumulates the
phase of $\lambda(\word{aa})$ over multiple traversals.

\subsection{The PMNS solar angle from the CKM manifold}
\label{sec:theta12nu}

The length-5 word $\word{abbAB}$ on $M_{\rm CKM}=\mfld{m006}$ gives
\begin{equation}
  180\degr - \phi(\word{abbAB}) = 180\degr - 146.380\degr = 33.620\degr
  \quad\text{vs.}\quad \theta_{12}^\nu=33.41\degr \quad (0.63\%).
\end{equation}
Notably, the solar neutrino mixing angle appears in the \emph{CKM manifold},
not the PMNS manifold.
This cross-manifold encoding suggests a geometric entanglement between the
quark and lepton sectors that goes beyond the simple K/N Iwasawa split
of our companion papers.

\subsection{The generator twist and CKM isospectrality}
\label{sec:generator}

The length-1 generator $\word{b}$ on $\mfld{m006}$ has
$\phi(\word{b})=89.16\degr\approx\pi/2$.
This near-right-angle twist propagates into all products: the three optimal
CKM words $\word{aaB}$, $\word{AbA}$, $\word{AAb}$ all give
$\phi\approx 92.49\degr$.
The triple isospectrality forces $J_{\rm CKM}\approx 0$ via phase cancellation,
explaining geometrically why the Jarlskog invariant is suppressed in the
quark sector~\cite{GentryCKM}.

\subsection{Full coincidence table}
\label{sec:full_table}

Table~\ref{tab:full} presents all angle matches within $5\degr$ and ratio
matches within $2\%$ from the complete census.
Each row is the canonical conjugacy class representative.

\begin{longtable}{llCCllr}
\caption{Complete A-factor spectral coincidences up to word length~5.
  Canonical conjugacy class representatives only.
  Fold: d$=$direct, $\pi$$=180\degr-|\phi|$, r$=$modulus ratio.}
\label{tab:full} \\
\toprule
Mfld & Word & \ell & |\phi|\degr & Fold & SM target & Error \\
\midrule
\endfirsthead
\toprule
Mfld & Word & \ell & |\phi|\degr & Fold & SM target & Error \\
\midrule
\endhead
\midrule\multicolumn{7}{r}{\small continued\ldots}\\
\endfoot
\bottomrule
\endlastfoot
%% ── 7xtheta12_CKM = 91.28 deg ──────────────────────────────────────────────
\multicolumn{7}{l}{\textit{Target: $7\theta_{12}^{\rm CKM}=91.28\degr$}} \\
\mfld{m003} & \word{AAAB} & 4 & 90.750 & d & $7\theta_{12}^{\rm CKM}$ & $0.530\degr$ (0.58\%) \\
\mfld{m006} & \word{B} & 1 & 90.844 & $\pi$ & $7\theta_{12}^{\rm CKM}$ & $0.436\degr$ (0.48\%) \\
\mfld{m006} & \word{AAb} & 3 & 92.487 & d & $7\theta_{12}^{\rm CKM}$ & $1.207\degr$ (1.32\%) \\
\mfld{m006} & \word{BBB} & 3 & 92.531 & d & $7\theta_{12}^{\rm CKM}$ & $1.251\degr$ (1.37\%) \\
\midrule
%% ── delta_CKM = 68.0 deg ────────────────────────────────────────────────────
\multicolumn{7}{l}{\textit{Target: $\delta_{\rm CKM}=68.0\degr$}} \\
\mfld{m003} & \word{BBBB} & 4 & 67.428 & d & $\delta_{\rm CKM}$ & $0.572\degr$ (0.84\%) \\
\mfld{m006} & \word{aa} & 2 & 67.654 & $\pi$ & $\delta_{\rm CKM}$ & $0.346\degr$ (0.51\%) \\
\mfld{m006} & \word{AAbbb} & 5 & 67.654 & d & $\delta_{\rm CKM}$ & $0.346\degr$ (0.51\%) \\
\mfld{m006} & \word{Ab} & 2 & 69.754 & d & $\delta_{\rm CKM}$ & $1.754\degr$ (2.58\%) \\
\midrule
%% ── theta12_NU = 33.41 deg ──────────────────────────────────────────────────
\multicolumn{7}{l}{\textit{Target: $\theta_{12}^\nu=33.41\degr$}} \\
\mfld{m003} & \word{ABaB} & 4 & 32.980 & d & $\theta_{12}^\nu$ & $0.430\degr$ (1.29\%) \\
\mfld{m006} & \word{abbAB} & 5 & 33.620 & $\pi$ & $\theta_{12}^\nu$ & $0.210\degr$ (0.63\%) \\
\mfld{m006} & \word{AAAbb} & 5 & 31.909 & d & $\theta_{12}^\nu$ & $1.501\degr$ (4.49\%) \\
\midrule
%% ── theta23_NU = 49.1 deg ───────────────────────────────────────────────────
\multicolumn{7}{l}{\textit{Target: $\theta_{23}^\nu=49.1\degr$}} \\
\mfld{m003} & \word{aabb} & 4 & 48.081 & $\pi$ & $\theta_{23}^\nu$ & $1.019\degr$ (2.1\%) \\
\mfld{m006} & \word{ABBAb} & 5 & 51.431 & d & $\theta_{23}^\nu$ & $2.331\degr$ (4.75\%) \\
\midrule
%% ── theta12_CKM = 13.04 deg ─────────────────────────────────────────────────
\multicolumn{7}{l}{\textit{Target: $\theta_{12}^{\rm CKM}=13.04\degr$}} \\
\mfld{m003} & \word{AAB} & 3 & 12.638 & $\pi$ & $\theta_{12}^{\rm CKM}$ & $0.402\degr$ (3.1\%) \\
\mfld{m003} & \word{AAAAb} & 5 & 13.535 & d & $\theta_{12}^{\rm CKM}$ & $0.495\degr$ (3.8\%) \\
\midrule
%% ── mb/mc = 3.2913 ──────────────────────────────────────────────────────────
\multicolumn{7}{l}{\textit{Target: $m_b/m_c=3.2913$}} \\
\mfld{m003} & \word{bbbb/bAbA} & 4 & --- & r & $m_b/m_c$ & $0.0003$ (0.01\%) \\
\mfld{m003} & \word{bbbb/AA} & 4 & --- & r & $m_b/m_c$ & $0.0003$ (0.01\%) \\
\mfld{m003} & \word{ABAbb/AB} & 5 & --- & r & $m_b/m_c$ & $0.0037$ (0.11\%) \\
\mfld{m003} & \word{AAbbb/ABB} & 5 & --- & r & $m_b/m_c$ & $0.0114$ (0.35\%) \\
\mfld{m006} & \word{AAbab/ABB} & 5 & --- & r & $m_b/m_c$ & $0.0190$ (0.58\%) \\
\mfld{m006} & \word{AAbab/Ab} & 5 & --- & r & $m_b/m_c$ & $0.0190$ (0.58\%) \\
\midrule
%% ── MZ/MW = 1.1345 ──────────────────────────────────────────────────────────
\multicolumn{7}{l}{\textit{Target: $M_Z/M_W=1.1345$}} \\
\mfld{m003} & \word{ABBBB/AAAB} & 5 & --- & r & $M_Z/M_W$ & $0.0011$ (0.10\%) \\
\mfld{m003} & \word{AAB/AAb} & 3 & --- & r & $M_Z/M_W$ & $0.0057$ (0.50\%) \\
\mfld{m003} & \word{ABB/A} & 3 & --- & r & $M_Z/M_W$ & $0.0082$ (0.72\%) \\
\mfld{m006} & \word{aaabb/baaab} & 5 & --- & r & $M_Z/M_W$ & $0.0010$ (0.09\%) \\
\mfld{m006} & \word{ABBB/B} & 4 & --- & r & $M_Z/M_W$ & $0.0053$ (0.46\%) \\
\mfld{m006} & \word{Abb/B} & 3 & --- & r & $M_Z/M_W$ & $0.0053$ (0.46\%) \\
\mfld{m006} & \word{ABBab/BBB} & 5 & --- & r & $M_Z/M_W$ & $0.0028$ (0.25\%) \\
\end{longtable}

%% ── 5. NULL RESULT ───────────────────────────────────────────────────────────
\section{Null Result: \texorpdfstring{$\theta_{13}^{\rm CKM}$}{theta13 CKM}}
\label{sec:null}

Six of seven independent SM flavor parameters are matched to better than $3\%$.
The sole outlier is $\theta_{13}^{\rm CKM}=0.201\degr$.
The minimum $|\phi|$ at length~5 on $\mfld{m006}$ is $\phi(\word{abbAb})=1.687\degr$,
giving a discrepancy of $1.49\degr$ ($740\%$).

The angle $\theta_{13}^{\rm CKM}$ is parametrically of order $\lambda_W^3\approx 0.011$
in the Wolfenstein parameterisation.
The spectral floor decreases approximately as
$\phi_{\rm min}(\ell)\sim\phi(\word{b})^{\ell/2}\approx(89.16\degr)^{\ell/2}$
under repeated composition of the near-$\pi/2$ generator; this estimate gives
$\phi_{\rm min}(12)\approx 0.18\degr$, placing $\theta_{13}^{\rm CKM}$ within
reach at word length $\ell^*\approx 12$--$14$.
This constitutes a falsifiable prediction of the framework.

%% ── 6. SELBERG ZETA AND SPECTRAL INTERPRETATION ─────────────────────────────
\section{Selberg Zeta Function and Spectral Interpretation}
\label{sec:selberg}

\subsection{First Laplace eigenvalues}

We compute $\lambda_1$ for both manifolds numerically via the Selberg zeta
function~\cite{Selberg56}
\begin{equation}
  Z_M(s) = \prod_{\gamma\,\rm prim}\prod_{m,n\geq 0}
  \Bigl(1-e^{i(m-n)\phi(\gamma)}e^{-(s+m+n)\ell(\gamma)}\Bigr),
  \label{eq:selberg}
\end{equation}
whose zeros on $\mathrm{Re}(s)=1$ correspond to Laplace eigenvalues via
$\lambda = s(2-s)$.
Using the first 200 primitive geodesics (length spectrum up to $\ell=5$) and
scanning the critical line, we find:
\begin{align}
  \lambda_1(\mfld{m003}) &= 2.480 \quad (r_1=1.493), \label{eq:lam1_m003}\\
  \lambda_1(\mfld{m006}) &= 2.821 \quad (r_1=1.604). \label{eq:lam1_m006}
\end{align}
Both values satisfy $\lambda_1 \geq 1/4$, consistent with the Selberg eigenvalue
conjecture, and lie well above the threshold for exceptional eigenvalues.

\subsection{Large spectral gap and SM-scale flavor parameters}

The values in~\eqref{eq:lam1_m003}--\eqref{eq:lam1_m006} are large compared
to generic hyperbolic 3-manifolds, where $\lambda_1$ can be made arbitrarily
small by Dehn surgery.
We propose the following heuristic interpretation.
The mixing angles of the SM satisfy
$\sin^2\theta_{ij}\lesssim 0.3$, which corresponds geometrically to
geodesic axis separations $\theta_{ij}\lesssim 35\degr$ in the Poincar\'e
ball picture of our CKM paper~\cite{GentryCKM}.
Manifolds with small $\lambda_1$ have long, nearly-flat geodesics whose
axes cluster at small separations, producing mixing angles much smaller or
larger than the SM values.
The large spectral gap $\lambda_1\sim 2.5$--$2.8$ restricts the geodesic
geometry to a ``sweet spot'' that naturally produces SM-scale mixing angles.
We conjecture that the condition $\lambda_1\gtrsim 2$ (equivalently,
$r_1\gtrsim 1.4$) is a necessary condition for a hyperbolic 3-manifold to
realise SM-scale mixing.

\subsection{Twist angles in the Selberg Euler factors}

Equation~\eqref{eq:selberg} shows that the twist angles $\phi(\gamma)$ appear
as phases in every Euler factor of $Z_M(s)$.
The spectral coincidences of Section~\ref{sec:results} therefore have a direct
spectral-theoretic interpretation: the SM flavor parameters are encoded in the
phases of the spectral determinant of the Laplace--Beltrami operator on $M$.
Concretely, the zeros of $Z_M(s)$ (and hence the heat kernel of $\Delta_M$)
are sensitive to the twist angles in a way that encodes the fermion flavor
structure.

\subsection{Odd-torsion conjecture and the Selberg zeta function}

For manifolds with $H_1(M;\mathbb{Z})$ containing 2-torsion, the holonomy
representation has real-valued characters on the torsion elements, introducing
additional real poles in $Z_M(s)$ that interfere constructively with the
principal series.
We conjecture that this interference is destructive for SM-scale angle
encoding: only manifolds with odd-prime torsion avoid this interference and
can realise the observed flavor hierarchy.
A rigorous version of this conjecture requires a trace formula computation
beyond the scope of this paper and is deferred to future work.

%% ── 7. DISCUSSION ────────────────────────────────────────────────────────────
\section{Discussion and Outlook}
\label{sec:discussion}

We have shown that the A-factor twist angle spectra of $\mfld{m003}$ and
$\mfld{m006}$ encode the SM flavor parameters with no free parameters.
The central results are:

\begin{enumerate}
\item $m_b/m_c$ to $0.01\%$ from a ratio of geodesic moduli on $\mfld{m003}$,
      connecting fermion masses to the real part of the complex length spectrum.
\item $\delta_{\rm CKM}=68.0\degr$ to $0.5\%$ from the length-2 generator word
      $\word{aa}$ on $\mfld{m006}$, appearing independently also at length 5.
\item The PMNS solar angle $\theta_{12}^\nu=33.41\degr$ to $0.6\%$ from a
      length-5 word on the \emph{CKM} manifold, suggesting cross-sector
      geometric entanglement.
\item $\lambda_1(\mfld{m003})=2.480$, $\lambda_1(\mfld{m006})=2.821$,
      both consistent with the Selberg conjecture and indicating a large
      spectral gap that we propose is geometrically responsible for
      SM-scale mixing angles.
\item $M_Z/M_W=1.1345$ to $0.09\%$ from a ratio of geodesic moduli on
      $\mfld{m006}$, a surprising result connecting electroweak symmetry
      breaking to hyperbolic geometry.
\item The sole null result, $\theta_{13}^{\rm CKM}=0.201\degr$, is predicted
      to appear at word length $\ell^*\approx 12$--$14$.
\end{enumerate}

Together with the companion papers, these results support the interpretation
of $\mfld{m003}$ and $\mfld{m006}$ as the geometrically natural homes of SM
flavor structure, with the Iwasawa decomposition $KAN$ distributing the
structure across mixing angles (K), CP phases and mass ratios (A), and
Borel structure (N).

Future directions include: (i) extending the census to length 7 to approach
$\theta_{23}^{\rm CKM}$ and test the $\theta_{13}^{\rm CKM}$ prediction;
(ii) a rigorous proof of the odd-torsion conjecture via arithmetic hyperbolic
geometry; (iii) investigating whether the $M_Z/M_W$ ratio has a natural
explanation in terms of a Kaluza--Klein reduction on $M$.

\section*{Acknowledgements}
The author thanks the \textsc{SnapPy} development team.
Computations used \textsc{SnapPy}~\cite{SnapPy}, \textsc{SciPy}, NumPy,
and Pandas.

\begin{thebibliography}{99}
\bibitem{PDG2024}
  S.~Navas \textit{et al.}\ (Particle Data Group),
  Phys.\ Rev.\ D \textbf{110}, 030001 (2024).
\bibitem{GentryCKM}
  M.~L.~Gentry,
  ``CKM Quark Mixing from Geodesic Axes of Compact Hyperbolic 3-Manifolds,''
  Phys.\ Rev.\ D, submitted March 2026 (es2026mar11\_966).
\bibitem{GentryPMNS}
  M.~L.~Gentry,
  ``PMNS Lepton Mixing from the Borel Structure of Hyperbolic Holonomy,''
  Phys.\ Rev.\ D, submitted March 2026 (es2026mar13\_942).
\bibitem{GentryCP}
  M.~L.~Gentry,
  ``CP Violation from the A-Factor of Hyperbolic Holonomy,''
  Phys.\ Rev.\ D, submitted March 2026 (es2026mar14\_722).
\bibitem{SnapPy}
  M.~Culler, N.~Dunfield, M.~Goerner, J.~Weeks,
  \textit{SnapPy}, a computer program for studying the geometry and topology
  of 3-manifolds, \url{http://snappy.computop.org}.
\bibitem{Selberg56}
  A.~Selberg,
  ``Harmonic analysis and discontinuous groups in weakly symmetric Riemannian
  spaces with applications to Dirichlet series,''
  J.\ Indian Math.\ Soc.\ \textbf{20}, 47--87 (1956).

\bibitem{GentryNPB1}
  M.~L.~Gentry,
  ``Geometric Origin of CP Phases from Hyperbolic Holonomy,''
  Nucl.\ Phys.\ B, submitted March 2026 (NPB-S-26-00539).
\bibitem{GentryNPB2}
  M.~L.~Gentry,
  ``Discrete Mixing Operators from Boundary Sector Geometry,''
  Nucl.\ Phys.\ B, submitted March 2026 (NPB-S-26-00540).
\bibitem{GentryNPB3}
  M.~L.~Gentry,
  ``Scale-Free Quadratic Forms, Symmetric Space Geometry, and Arithmetic
  Logarithmic Lattices,''
  Nucl.\ Phys.\ B, submitted March 2026 (NPB-S-26-00538).
\end{thebibliography}

\end{document}



